% =====================================================================
% ejercicio02.tex
% Responsabilidad única: plantear y resolver el sistema normal
% X^T X beta = X^T y, interpretar coeficientes y discutir colinealidad
% perfecta y el efecto de añadir variables sobre el error de
% entrenamiento y de generalización. Cifras verificadas con
% verificacion/teoria.py.
% =====================================================================
\subsection{Ejercicio 2 --- Regresión lineal múltiple: solución matricial y colinealidad}

Se modela la demanda $y$ (centenas de unidades) como
$\hat y=\beta_0+\beta_1x_1+\beta_2x_2$, donde $x_1$ es la inversión
publicitaria (unidades de presupuesto) y $x_2\in\{0,1\}$ indica si hubo
promoción.

\begin{enumerate}

\item \textbf{Matriz $X$, $X^\top X$ y $X^\top y$.}
Agregando la columna de unos del intercepto,
\[
    X = \begin{pmatrix}
        1 & 0 & 0\\ 1 & 1 & 0\\ 1 & 2 & 0\\ 1 & 0 & 1\\ 1 & 1 & 1\\ 1 & 2 & 1
    \end{pmatrix}\in\mathbb{R}^{6\times 3},
    \qquad
    y = \begin{pmatrix} 2\\4\\5\\5\\6\\8 \end{pmatrix}\in\mathbb{R}^{6\times 1},
    \qquad
    \beta\in\mathbb{R}^{3\times1},
\]
y cada entrada de $X^\top X$ es un producto punto entre columnas de $X$
(por ejemplo $\sum_i x_{i1}^2=0+1+4+0+1+4=10$ y
$\sum_i x_{i1}x_{i2}=0+1+2=3$), así como cada entrada de $X^\top y$ es
el producto punto de una columna con $y$, de forma que
\[
    X^\top X = \begin{pmatrix}
        n & \sum x_{1} & \sum x_{2}\\
        \sum x_{1} & \sum x_{1}^2 & \sum x_{1}x_{2}\\
        \sum x_{2} & \sum x_{1}x_{2} & \sum x_{2}^2
    \end{pmatrix}
    = \begin{pmatrix} 6 & 6 & 3\\ 6 & 10 & 3\\ 3 & 3 & 3 \end{pmatrix},
    \qquad
    X^\top y = \begin{pmatrix} \sum y\\ \sum x_1y\\ \sum x_2y \end{pmatrix}
    = \begin{pmatrix} 30\\ 36\\ 19 \end{pmatrix}.
\]

\item \textbf{Derivación y solución del sistema normal.}
Expandiendo la función objetivo,
\[
    S(\beta) = (y-X\beta)^\top(y-X\beta)
    = y^\top y - 2\beta^\top X^\top y + \beta^\top X^\top X\beta ,
\]
y usando las reglas $\nabla_\beta(\beta^\top a)=a$ y
$\nabla_\beta(\beta^\top A\beta)=2A\beta$ para $A$ simétrica, se obtiene
\[
    \nabla_\beta S = -2X^\top y + 2X^\top X\beta = 0
    \;\Longrightarrow\;
    X^\top X\hat\beta = X^\top y ,
\]
que corresponde a un mínimo porque la hessiana $2X^\top X$ es semidefinida
positiva ($v^\top X^\top Xv=\lVert Xv\rVert^2\ge0$), y definida positiva
cuando las columnas de $X$ son linealmente independientes, como en este
caso ($\det X^\top X=36\neq0$). El sistema a resolver es
\[
    \begin{aligned}
        6\beta_0 + 6\beta_1 + 3\beta_2 &= 30 &\quad&(\mathrm{E}_1)\\
        6\beta_0 + 10\beta_1 + 3\beta_2 &= 36 &&(\mathrm{E}_2)\\
        3\beta_0 + 3\beta_1 + 3\beta_2 &= 19 &&(\mathrm{E}_3)
    \end{aligned}
\]
y se procede por eliminación: primeramente $\mathrm{E}_2-\mathrm{E}_1$
da $4\beta_1=6$, luego $\mathrm{E}_1-2\mathrm{E}_3$ da $-3\beta_2=-8$, y
finalmente al sustituir ambos en $\mathrm{E}_3$ se despeja $\beta_0$:
\[
    \hat\beta_1 = \frac{3}{2} = 1.5, \qquad
    \hat\beta_2 = \frac{8}{3} \approx 2.6667, \qquad
    3\hat\beta_0 = 19 - 4.5 - 8 = 6.5
    \;\Rightarrow\; \hat\beta_0 = \frac{13}{6} \approx 2.1667 .
\]

\item \textbf{Interpretación, predicción y MSE.}
\begin{itemize}
    \item $\hat\beta_1=1.5$: manteniendo fija la condición de promoción,
    cada unidad adicional de presupuesto publicitario se asocia con
    $1.5$ centenas ($150$ unidades) más de demanda predicha.
    \item $\hat\beta_2\approx2.6667$: para un mismo nivel de inversión
    publicitaria, realizar la promoción se asocia con unas $2.6667$
    centenas (aproximadamente $267$ unidades) más de demanda predicha.
\end{itemize}
Para $(x_1,x_2)=(1.5,1)$ se predice
$\hat y=\tfrac{13}{6}+1.5\cdot1.5+\tfrac83=\tfrac{85}{12}\approx7.0833$
centenas, es decir, unas $708$ unidades. Con las predicciones de
entrenamiento
\[
    \hat y\approx(2.1667,\;3.6667,\;5.1667,\;4.8333,\;6.3333,\;7.8333),
\]
los residuos son $e=\pm\tfrac16$ o $\pm\tfrac13$ y
\[
    MSE = \frac{1}{6}\Bigl(4\cdot\tfrac{1}{36} + 2\cdot\tfrac{1}{9}\Bigr)
    = \frac{1}{6}\cdot\frac{1}{3} = \frac{1}{18} \approx 0.0556 .
\]
Estos coeficientes no prueban un efecto causal, porque solo describen
la mejor combinación lineal \emph{en estos datos}, y los datos no
provienen (hasta donde se sabe) de un experimento donde la inversión o
la promoción se asignaran al azar, de forma que podría existir una
variable omitida (temporada, precio, competencia) que mueva a la vez la
decisión de invertir y la demanda; además con $n=6$ observaciones
didácticas no hay forma de separar el efecto real del ruido, así que la
lectura correcta es de \emph{asociación condicional}, no de
intervención.

\item \textbf{Colinealidad perfecta ($x_3 = 2x_1$).}
Si se agrega $x_3=2x_1$, la nueva matriz
$\tilde X=[\mathbf{1},x_1,x_2,x_3]$ satisface
$0\cdot\mathbf{1}+2\,x_1+0\cdot x_2-1\cdot x_3=\mathbf{0}$, es decir,
existe $v=(0,2,0,-1)^\top\neq0$ con $\tilde Xv=0$, por lo que sus
columnas son linealmente dependientes, $\operatorname{rango}(\tilde X)=3<4$
y $\tilde X^\top\tilde X$ (de $4\times4$) es singular. En consecuencia,
si $\hat\beta$ es una solución, también lo es $\hat\beta+c\,v$ para todo
$c\in\mathbb{R}$, porque $\tilde X(\hat\beta+cv)=\tilde X\hat\beta$:
cualquier combinación con $\beta_1+2\beta_3=1.5$ (y los mismos
$\beta_0,\beta_2$) produce \textbf{exactamente las mismas
predicciones}, de modo que el ajuste sigue siendo único en términos de
$\hat y$, pero no en términos de coeficientes, que dejan de ser
interpretables individualmente (no tiene sentido preguntar ``el efecto
de $x_1$ con $x_3$ fijo'' si $x_3$ no puede quedarse fijo mientras $x_1$
cambia). Para ajustar sin invertir una matriz singular se puede, en
orden de preferencia:
\begin{itemize}
    \item Eliminar la variable redundante: Se descarta $x_3$ (o $x_1$),
    lo que devuelve exactamente el modelo del inciso b.
    \item Usar la pseudoinversa de Moore--Penrose: Se calcula
    $\hat\beta=\tilde X^{+}y$ (mediante descomposición en valores
    singulares, \texttt{np.linalg.lstsq} o \texttt{pinv}), que elige
    entre las infinitas soluciones la de norma mínima, en este caso
    $\hat\beta\approx(2.1667,\,0.3,\,2.6667,\,0.6)$, con
    $0.3+2\cdot0.6=1.5$ y las mismas predicciones.
    \item Regularizar: Se resuelve
    $(\tilde X^\top\tilde X+\lambda I)\beta=\tilde X^\top y$ con
    $\lambda>0$ (regresión ridge), sistema que siempre es invertible.
\end{itemize}

\item \textbf{Más variables: error de entrenamiento vs.\ datos nuevos.}
Si se añade una variable $x_{p+1}$, el modelo ampliado contiene al
original como caso particular (basta fijar $\beta_{p+1}=0$), de forma
que el mínimo del SSE sobre un conjunto más grande de candidatos no
puede ser mayor:
\[
    \min_{\beta\in\mathbb{R}^{p+2}} \lVert y-\tilde X\beta\rVert^2
    \;\le\;
    \lVert y-\tilde X(\hat\beta_{\text{orig}},0)\rVert^2
    = \min_{\beta\in\mathbb{R}^{p+1}} \lVert y-X\beta\rVert^2 ,
\]
y por eso el error de entrenamiento nunca aumenta (en el peor caso se
mantiene igual, cuando la nueva variable no aporta nada en la muestra).
Sin embargo, en datos nuevos sí puede empeorar, porque el nuevo
coeficiente también se ajusta al ruido particular de la muestra de
entrenamiento, lo cual aumenta la varianza de las predicciones sin
reducir necesariamente el sesgo (sobreajuste), y con pocas
observaciones o predictores correlacionados este efecto se vuelve más
fuerte, razón por la cual en la Parte II la comparación entre modelos se
decide con un conjunto de validación y no con el error de
entrenamiento.

\end{enumerate}
